% ============================================
% Letter of Interest – Northwestern College
% ============================================

\documentclass[11pt,letterpaper]{letter}

\usepackage{geometry}
\usepackage{microtype}
\usepackage[utf8]{inputenc}

\geometry{left=25mm,right=25mm,top=25mm,bottom=25mm}

\signature{Decory Edwards}
\address{
Baltimore, MD \\
\texttt{dedwar65@jhu.edu}
}

\begin{document}

\begin{letter}{
Vice President for Academic Affairs \\
Northwestern College \\
Orange City, Iowa
}

\opening{Dear Members of the Search Committee,}

Thank you for the opportunity to complete my application for the Visiting Faculty position in Economics. I am currently a Ph.D.\ candidate in Economics at Johns Hopkins University, expecting to complete my degree in May 2026. I am grateful for the chance to describe my teaching experience, scholarly trajectory, and how Northwestern College’s Christian Identity, Vision for Learning, and Vision for Diversity would shape my teaching and scholarship.

\textbf{Teaching Experience}

My teaching experience has centered on undergraduate instruction in economics and quantitative methods. I approach teaching with clarity, structure, and enthusiasm, emphasizing economic reasoning as a disciplined way of thinking about stewardship, incentives, tradeoffs, and human flourishing. In principles-level courses, I prioritize accessibility while maintaining rigor, helping students develop strong analytical foundations. In upper-level courses, I integrate real data, empirical reasoning, and policy discussion to deepen students’ engagement with complex economic issues.

I value mentorship and see teaching as relational as well as intellectual. I strive to create a classroom environment where students feel supported, challenged, and encouraged to connect economic concepts to ethical and social questions. I believe economics education should cultivate careful reasoning, intellectual humility, and thoughtful engagement with the world.

\textbf{Scholarly Trajectory}

My research focuses on macroeconomics and household finance, particularly how institutional environments and differences in financial returns shape wealth accumulation over the life cycle. In my dissertation, I develop and calibrate heterogeneous-agent macroeconomic models to match empirical wealth data. I also conduct empirical work using large longitudinal datasets to examine how institutional trust and economic behavior interact over time.

Moving forward, I aim to continue research that bridges theory and applied analysis, particularly in areas related to economic opportunity, inequality, financial decision-making, and institutional design. I am committed to scholarship that maintains technical rigor while addressing questions that matter for communities and public life. In a liberal arts context, I see scholarship not only as publication, but as modeling intellectual curiosity and disciplined inquiry for students.

\textbf{Christian Identity, Vision for Learning, and Vision for Diversity}

Northwestern College’s Reformed Christian identity and its emphasis on integrating faith and learning resonate deeply with me. I believe that economics, when taught thoughtfully, can illuminate questions of stewardship, justice, responsibility, and service. The Vision for Learning’s emphasis on developing thoughtful, compassionate leaders aligns with my goal of helping students think critically about economic systems while recognizing their moral and social dimensions.

I also value Northwestern’s Vision for Diversity, which emphasizes unity in Christ alongside engagement across differences. In the classroom, I encourage respectful dialogue and thoughtful engagement with diverse perspectives, recognizing that economic issues often intersect with cultural, social, and ethical questions. I strive to cultivate intellectual hospitality—welcoming students from varied backgrounds and encouraging them to engage complex topics with humility and care.

Ultimately, I view teaching economics in a Christian liberal arts setting as an opportunity to help students develop disciplined analytical skills, moral imagination, and a commitment to service. I would be honored to contribute to Northwestern College’s mission of preparing students for lives of faithful leadership and meaningful impact.

Thank you for your consideration. I appreciate the opportunity to be considered for this role and look forward to the possibility of contributing to the Northwestern College community.

\closing{Sincerely,}

\end{letter}
\end{document}